\chapter{QUATTRO FOTOGRAFIE DEL DELTA DEL KIRIOLOJ}\label{quattro-fotografie-del-delta-del-kirioloj}

\hfill\break

È difficile catturare la brutale follia dei tempi.

Alcuni giorni sembrava quasi liberatoria. Al di là della nostra
insignificante illusione del cielo, il Sole continuava a espandersi, le
stelle si spegnevano o nascevano, un pianeta morto era stato infuso di
vita e aveva fatto evolvere una civiltà equivalente o superiore alla
nostra. Da noi, invece, i governi venivano fatti cadere e sostituiti e i
loro rimpiazzi rovesciati; le religioni, filosofie e ideologie si
trasformavano, fondevano e generavano stirpi mutanti. Il vecchio mondo
ordinato si stava sgretolando.

Dalle rovine però crescevano cose nuove. Coglievamo l'amore ancora verde
e assaporavamo la sua acerbità: presumo che Molly Seagram mi amasse
soprattutto perché ero disponibile. E perché no? L'estate stava finendo
e il raccolto era incerto.

Il movimento del Nuovo Regno, estinto da tempo, veniva considerato
lungimirante ma allo stesso tempo fortemente superato, la sua timida
ribellione contro l'antico consenso ecclesiastico era ombra di nuove
devozioni più estreme. Nel mondo occidentale sorsero ovunque culti
dionisiaci, spogliati della religiosità e dell'ipocrisia del vecchio NR:
club del sesso con bandiere e simboli sacri. Non disdegnavano la gelosia
umana, anzi la abbracciavano o addirittura ne godevano: gli amanti
disprezzati prediligevano le pistole calibro 45 a distanza ravvicinata,
una rosa rossa sul corpo della vittima. Era la Tribolazione ridelineata
in dramma elisabettiano.

Simon Townsend, se fosse nato dieci anni dopo, sarebbe potuto incappare
in uno di questi generi di spiritualità alla Quentin Tarantino. Ma il
fallimento dell'NR lo aveva deluso e desiderava ardentemente qualcosa di
più semplice. Diane mi chiamava ancora di tanto in tanto - più o meno
una volta al mese, sotto i giusti auspici e quando Simon non era in casa
- per aggiornarmi sulla sua situazione o semplicemente per rievocare il
passato, attizzando i ricordi come fossero braci per riscaldarsi al loro
calore. Calore che sembrava quasi assente in casa, malgrado la sua
situazione finanziaria fosse un po' migliorata.

Simon si occupava a tempo pieno della manutenzione del Jordan
Tabernacle, la loro chiesetta indipendente; Diane svolgeva mansioni
d'ufficio temporanee, un lavoro saltuario che spesso la costringeva ad
aggirarsi nervosa per l'appartamento senza far niente o le consentiva di
svignarsela nella biblioteca locale per leggere libri che Simon
disapprovava: romanzi contemporanei, attualità. Diceva che il Jordan
Tabernacle era una chiesa ``di rottura'', che i fedeli erano
incoraggiati a spegnere la TV ed evitare libri, giornali e altri
prodotti culturali effimeri. Altrimenti rischiavano di affrontare
l'Assunzione in cielo in una condizione impura.

Diane non difendeva mai quelle idee - non mi faceva mai la predica - ma
le accettava, facendo attenzione a non metterle in discussione. A volte
mi spazientivo un po'. «Diane,» le dissi una sera, «ci credi davvero a
queste cose?»

«A quali ``cose'', Tyler?»

«Scegli tu. Il fatto di non tenere libri in casa. Che gli Ipotetici sono
agenti della Parusia. Tutta quella merda.» (Forse avevo bevuto una birra
di troppo.)

«Simon ci crede.»

«Non ti ho chiesto cosa ne pensa Simon.»

«Simon è più devoto di me. Lo invidio. E so già la scena che hai in
mente, magari lui che mi grida: ``Getta quei libri nella spazzatura!''
Tu lo immagini come un mostro, un essere arrogante. Ma non lo è. Le sue
scelte sono veri atti di umiltà e sottomissione. Simon riesce a
concedersi a Dio in un modo a me sconosciuto.»

«Beato lui.»

«E lo è, davvero. Tu non capisci, ma lui è in pace. Al Jordan ha trovato
la serenità. Può guardare in faccia lo Spin e sorridergli, perché sa di
essere salvo.»

«E tu? Tu non sei salva?»

Lasciò scorrere nella linea telefonica un lungo silenzio.

«Bella domanda\ldots{} davvero. Continuo a pensare che forse la mia fede
non c'entri. Forse la fede di Simon è sufficiente per entrambi, forse è
abbastanza potente da trainarmi per un po'. In realtà è stato molto
paziente, con me. L'unica cosa su cui litighiamo è il fatto di avere
figli. Simon vorrebbe averne. La chiesa incoraggia ad averne. E lo
capisco, ma i soldi sono pochi, e poi\ldots{} con il mondo in queste
condizioni\ldots»

«Non è una decisione da prendere sotto pressione.»

«Non voglio insinuare che mi stia facendo pressioni. Mi dice:
``Rimettiti nelle mani di Dio. Rimettiti nelle mani di Dio e andrà tutto
bene.''»

«Ma tu sei troppo intelligente per credere a queste cose.»

«Davvero? Oh, Tyler, spero di no. Spero che non sia vero.»

\separatorefiori

Molly, invece, non sapeva che farsene di ``tutte queste stronzate su
Dio'', come le chiamava lei. ``Ognuna per sé'' era la filosofia di Moll.
A maggior ragione, diceva, visto che il mondo stava perdendo il
controllo e nessuno di noi sarebbe campato oltre i cinquant'anni. «Non
intendo passare inginocchiata il tempo che mi resta.»

Era tosta per natura.

I genitori di Molly erano allevatori di vacche. Avevano passato dieci
anni in mezzo a controversie legali per un progetto di estrazione del
petrolio da sabbie bituminose che confinavano con la loro proprietà e la
stavano lentamente avvelenando. Alla fine, con una transazione
stragiudiziale abbastanza cospicua, avevano ceduto il loro ranch in
cambio di una bella pensione per loro e un'istruzione decente per la
figlia. Ma era il genere di esperienza, diceva Molly, che avrebbe fatto
venire i calli al culo anche a un angelo.

Del panorama sociale che evolveva poche cose la sorprendevano. Una sera
ci sedemmo davanti alla TV a guardare un servizio sulle rivolte di
Stoccolma. Una folla di pescatori di merluzzo e radicali religiosi
scagliava mattoni contro le finestre e bruciava automobili; gli
elicotteri della polizia cosparsero sui dimostranti un gel viscoso
finché gran parte della Gamla Stan assunse l'aspetto di un grumo di
catarro sputato da un Godzilla tisico. Feci un'osservazione stupida su
quanto ci si possa comportare male quando si è spaventati, e Molly
disse: «Ma dai, Tyler, sul serio provi simpatia per questi stronzi?»

«Non ho detto questo, Moll.»

«Con la scusa dello Spin, hanno carta bianca per distruggere la sede del
loro Parlamento? Perché, perché sono spaventati?»

«Non è una scusa. È un motivo. Non hanno un futuro. Credono di essere
condannati.»

«Condannati a morire. Be', benvenuti nella condizione umana. Moriranno
loro, morirai tu, morirò io\ldots{} e quando mai non è stato così?»

«Siamo tutti mortali, ma almeno prima avevamo la consolazione di sapere
che la specie umana sarebbe proseguita senza di noi.»

«Ma anche le specie sono mortali. L'unica cosa che è cambiata è che
all'improvviso non c'è via d'uscita da un futuro nebuloso. Può darsi che
tra qualche anno moriremo tutti insieme in un modo spettacolare\ldots{}
ma anche questa è ancora soltanto una possibilità. Gli Ipotetici
potrebbero mantenerci in vita più a lungo. Per una qualunque ragione a
noi incomprensibile.»

«E questo non ti spaventa?»

«Ma certo! Tutto questo mi spaventa. Ma non è un motivo per scendere in
strada ad ammazzare gente.» Indicò la TV. Qualcuno aveva lanciato una
granata nel Riksdag. «È una cosa così\ldots{} stupida. Non si ottiene
nulla, in questo modo. È una prova di forza dettata dagli ormoni. Roba
da scimmie.»

«Non puoi far finta che la cosa non ti tocchi.»

La sua risata mi colse di sorpresa. «No\ldots{} questo è il \emph{tuo}
stile, non il mio.»

«Ah, sì?»

Abbassò la testa ma la rialzò fissandomi, quasi in segno di sfida. «Sì.
È quello che fai con lo Spin, fingendo che non sia un problema. E fai lo
stesso con i Lawton. Loro ti usano, ti ignorano, e tu sorridi come fosse
l'ordine naturale delle cose.»

Mi osservava come se si aspettasse una mia reazione. Ero troppo testardo
per dargliela vinta e lei continuò: «Penso solo che ci siano modi
migliori per trascorrere il tempo che manca alla fine del mondo.»

Ma non disse quali fossero questi modi migliori.

\separatorefiori

Al momento dell'assunzione, tutti quelli che lavoravano alla Perielio,
me compreso, avevano firmato un accordo di riservatezza e tutti noi
eravamo stati sottoposti a un controllo accurato delle referenze e a un
esame per il rispetto delle norme di sicurezza nazionale. Eravamo
discreti e rispettavamo la necessità di mantenere all'interno
dell'azienda le conversazioni confidenziali. Le fughe di notizie
avrebbero potuto spaventare i comitati del Congresso, imbarazzare amici
potenti, scoraggiare finanziatori.

Ma adesso c'era un marziano che viveva nel campus - quasi tutta l'ala
nord era stata adibita ad alloggio temporaneo per Wun Ngo Wen e i suoi
addetti - e quel segreto era difficile da mantenere.

In ogni caso non lo si sarebbe potuto tenere nascosto ancora per molto.
Nel momento in cui Wun arrivò in Florida, gran parte dell'élite di
Washington e diversi capi di Stato stranieri ne avevano già sentito
parlare. Il Dipartimento di Stato gli aveva concesso uno status
giuridico ad hoc e prevedeva di presentarlo a tempo debito a livello
internazionale. I suoi addetti lo stavano preparando per l'inevitabile
pressione mediatica.

Avrebbero potuto e forse dovuto gestire il suo arrivo in modo diverso.
Avrebbero potuto presentarlo le Nazioni Unite, rendere pubblica fin da
subito la sua presenza. L'amministrazione Garland avrebbe sicuramente
ricevuto forti critiche per averlo tenuto nascosto. Il Partito Cristiano
Conservatore già iniziava con le insinuazioni: ``Il governo sa più di
quanto dice sui risultati del progetto di terraformazione'' con la
speranza di far uscire allo scoperto il Presidente o gettare Lomax, suo
aspirante successore, in pasto alle critiche. Critiche che
inevitabilmente ci sarebbero state; ma Wun aveva espresso il desiderio
di non diventare argomento di campagna elettorale. Voleva mostrarsi al
pubblico ma per rivelarsi, disse, avrebbe aspettato fino a novembre.

L'esistenza di Wun Ngo Wen, comunque, era solo il segreto più evidente
tra quelli che riguardavano il suo arrivo. Ce n'erano altri. Si
prospettava una strana estate, alla Perielio.

Quell'agosto Jason mi convocò all'ala nord. Lo incontrai nel suo ufficio
- il suo vero ufficio, non la suite arredata con gusto dove riceveva i
visitatori ufficiali e la stampa: un cubo senza finestre con una
scrivania e un divano. Seduto sul bordo della sedia tra pile di riviste
scientifiche, in Levi's e felpa unta, sembrava spuntato dal disordine
come una pianta idroponica. Sudava. Mai un buon segno, con Jase.

«Sto perdendo di nuovo l'uso delle gambe» disse.

Liberai uno spazio sul divano, mi sedetti e aspettai che mi desse altri
particolari.

«Da un paio di settimane ci sono brevi episodi. La solita cosa,
formicolio al mattino. Niente che non possa sopportare. Ma non va via.
In realtà sta peggiorando. Penso che dovremmo cambiare terapia.»

Forse sì. Ma non mi piaceva per niente quello che gli stava causando la
cura. Jase ormai prendeva ogni giorno una manciata di pillole:
potenziatori della mielina per rallentare la perdita del tessuto
nervoso, richiami neurologici per aiutare il cervello a riavviare zone
danneggiate e farmaci secondari per trattare gli effetti collaterali
della cura principale. Potevamo aumentargli il dosaggio? Forse. Ma il
procedimento aveva un tetto di tossicità ed eravamo già al limite. Aveva
perso peso e forse anche qualcosa di più importante: un certo equilibrio
emotivo. Jase parlava più velocemente del solito e sorrideva meno
spesso. Pochi anni prima era del tutto a proprio agio nel suo corpo,
adesso invece si muoveva come una marionetta - quando allungò il braccio
per afferrare una tazza, la mano superò il bersaglio e poi tornò
indietro per riuscire a prenderla.

«In ogni caso,» dissi, «dobbiamo sentire il parere del dottor
Malmstein.»

«Non mi è davvero possibile uscire da qui per incontrarlo. Le cose sono
cambiate, se non te ne sei accorto. Non possiamo consultarlo al
telefono?»

«Forse. Chiederò.»

«E, nel frattempo, puoi farmi un altro favore?»

«Di che si tratta, Jase?»

«Spiega il mio problema a Wun. Procuragli un paio di manuali
sull'argomento.»

«Testi di medicina? Perché, è un medico?»

«Non proprio, ma ha portato con sé molte informazioni. Le scienze
biologiche marziane sono notevolmente in anticipo rispetto alle nostre.»
(Lo disse con un sorriso sbilenco che non riuscii a interpretare.)
«Crede di riuscire ad aiutarci.»

«Sei serio?»

«Serissimo. Smettila di guardarmi scioccato. Ne vuoi parlare con lui?»

Un uomo di un altro pianeta. Un uomo con centomila anni di storia
marziana alle spalle. «Be', sì. Ne sarei onorato. Ma\ldots»

«Allora organizzo la cosa.»

«Ma se ha quel genere di conoscenza medica che può curare in modo
efficace la SMA, deve mettersi in contatto con medici più bravi di me.»

«Wun ha portato con sé intere enciclopedie. C'è già chi sta esaminando
gli archivi marziani - alcune parti, almeno - per cercare informazioni
utili, mediche e non. Questo è solo un evento secondario.»

«Sono sorpreso che possa perdere tempo per un evento secondario.»

«Si annoia più spesso di quanto tu creda. E poi non ha molti amici.
Credo che gli farebbe piacere passare un po' di tempo con qualcuno che
non lo ritenga un salvatore o una minaccia. Comunque, vorrei che nei
prossimi giorni parlassi con Malmstein.»

«Certo.»

«E chiamalo da casa tua, d'accordo? Non mi fido più dei telefoni, qui
dentro.»

Sorrise come se avesse detto qualcosa di divertente.

\separatorefiori

Di tanto in tanto quell'estate facevo delle passeggiate sulla spiaggia
pubblica oltre la strada di fronte al mio appartamento.

Non aveva molto della spiaggia. Una lunga lingua di terra non edificata
la proteggeva dall'erosione e la rendeva inutile per i surfisti. Nei
pomeriggi caldi i vecchi motel scrutavano la sabbia con occhi di vetro e
qualche turista mogio si bagnava i piedi nella spuma delle onde.

Scesi e mi sedetti su una passerella di legno rovente sospesa su una
macchia di cespugli, a osservare le nubi addensarsi all'orizzonte verso
est e a pensare a ciò che aveva detto Molly, che fingevo che lo Spin non
fosse un problema (e anche i Lawton), simulando una serenità d'animo che
di fatto non potevo possedere.

Volevo dare a Molly ciò che le spettava. Credo fosse il modo in cui
tentavo di prendermi cura di lei.

``Spin'' era un nome stupido ma inevitabile per quello che era stato
fatto alla Terra. Cioè, non era una descrizione corretta della fisica -
niente in realtà girava più forte o più veloce di prima - ma era una
metafora calzante. In realtà la Terra era più statica di quanto non
fosse mai stata, anche se \emph{sembrava} che ruotasse fuori controllo.
Almeno per gli aspetti importanti. Per questo avevamo bisogno di
aggrapparci a qualcosa, a meno di non volere scivolare nell'oblio.

Quindi forse mi stavo aggrappando ai Lawton - non solo a Jason e Diane
ma a tutto il loro mondo, la Grande Casa e la Piccola Casa, per fedeltà
all'infanzia perduta. Forse quello era l'unico appiglio a cui potevo
aggrapparmi. E forse non era necessariamente un male. Se Moll aveva
ragione, tutti noi dovevamo aggrapparci a qualcosa o saremmo stati
perduti. Diane si era aggrappata alla fede, Jason si era aggrappato alla
scienza.

E io mi ero aggrappato a Jason e Diane.

Me ne andai dalla spiaggia quando comparvero le nuvole, preannunciando
uno di quegli inevitabili temporali pomeridiani di fine agosto, il cielo
a est inquieto e lampeggiante di fulmini, la pioggia che iniziava a
sferzare i balconi tristi e color pastello dei motel. Il tempo di
tornare a casa, e avevo già i vestiti bagnati. Ci vollero ore per
asciugarli nell'aria umida. La tempesta passò all'imbrunire, ma lasciò
una quiete fetida e umida.

Molly passò dopo cena e scaricammo un film contemporaneo, uno di quei
drammi da salotto vittoriano di cui lei era appassionata. Dopo il film
andò in cucina per prepararci qualcosa da bere mentre io chiamai David
Malmstein dal telefono nella stanza degli ospiti. Malmstein disse che
avrebbe visitato Jase non appena gli fosse stato possibile, ma pensava
che sarebbe stato opportuno aumentare un po' la terapia, purché sia io
che Jase prestassimo attenzione a qualsiasi reazione spiacevole.

Riattaccai il telefono, lasciai la stanza e trovai Molly in sala con due
bicchieri in mano e un'espressione perplessa sul volto. «Dove sei
andato?»

«Solo a fare una chiamata.»

«Qualcosa di importante?»

«No.»

«Un controllo su un paziente?»

«Qualcosa di simile.»

\separatorefiori

Dopo qualche giorno Jase organizzò un incontro tra me e Wun Ngo Wen
nell'alloggio di Wun alla Perielio.

L'ambasciatore marziano viveva in una stanza che aveva arredato da
catalogo secondo il proprio gusto. I mobili erano leggeri, bassi e in
vimini. Un tappeto di stoffa riciclata ricopriva il pavimento in
linoleum. Su un semplice tavolo di pino grezzo stava poggiato un
computer. C'erano librerie coordinate alla scrivania. A quanto pareva i
marziani decoravano le case come le matricole dei college.

Consegnai a Wun il materiale tecnico che voleva: un paio di libri
sull'eziologia e sul trattamento della sclerosi multipla, oltre a una
serie di articoli del JAMA sulla SMA. La SMA, nel pensiero corrente, non
era affatto la SM; era una malattia completamente diversa, un disturbo
genetico con sintomi analoghi alla SM e un invecchiamento simile delle
guaine mieliniche che proteggono il tessuto nervoso umano. La SMA si
distingueva per la gravità, la rapida progressione e la resistenza alle
terapie standard.

Wun mi disse che non aveva familiarità con la patologia ma avrebbe
cercato informazioni nei suoi archivi.

Lo ringraziai ma sollevai un'ovvia obiezione: non era un dottore, e la
fisiologia marziana era parecchio insolita - se anche avesse trovato una
terapia idonea, avrebbe funzionato nel caso di Jason?

«Non siamo così diversi come sembrerebbe. Una delle prime cose che la
tua gente ha fatto è stata quella di sequenziare il mio genoma. È
indistinguibile dal vostro.»

«Non volevo offenderti.»

«Non mi sono offeso. Centomila anni sono un distacco lungo, tanto lungo
da provocare quello che i biologi chiamano un evento di speciazione.
Guarda caso, però, la tua gente e la mia sono pienamente interfertili.
Le differenze evidenti tra di noi sono adattamenti superficiali a un
ambiente più fresco e più asciutto.»

Parlava con un'autorevolezza che superava la sua statura. Aveva la voce
più acuta della media degli adulti, sebbene non giovanile; era ritmata,
quasi femminile, ma sempre da uomo di Stato.

«Comunque,» dissi, «ci sono potenziali problemi legali se parliamo di
una terapia che non è passata per il processo di approvazione della
FDA.»

«Sono sicuro che Jason sarebbe disposto ad aspettare l'approvazione
ufficiale ma non credo che la sua malattia sia così paziente.» E a quel
punto Wun sollevò la mano per anticipare ulteriori obiezioni. «Lasciami
leggere quello che mi hai portato. Poi ne ridiscuteremo.»

Conclusa la questione più urgente, mi chiese di restare a parlare. Ne
ero lusingato. Malgrado la sua stranezza, c'era qualcosa di confortante
nella presenza di Wun, come se emanasse un'aura di calma. Si mise a
sedere sulla voluminosa sedia di vimini, i piedi penzoloni, e ascoltò
con vivo interesse un breve accenno sulla mia vita. Fece un paio di
domande su Diane (``Jason non parla molto della sua famiglia'') e altre
sulla scuola di medicina (il concetto di dissezionare i cadaveri era
nuovo per lui; sobbalzò quando glielo descrissi\ldots{} quasi tutti
hanno la stessa reazione).

E quando gli chiesi della sua vita infilò la mano nella piccola
cartellina grigia che teneva con sé e mi mostrò una serie di immagini
stampate, fotografie che aveva portato in formato digitale. Quattro foto
di Marte.

«Solo quattro?»

Scrollò le spalle. «Nessun numero, per quanto grande, può sostituire i
ricordi. E naturalmente negli archivi ufficiali c'è molto più materiale
visivo. Queste sono mie. Personali. Ti interessa vederle?»

«Sì, certo.»

Me le passò.

\emph{Foto 1}: Una casa. Era evidentemente un luogo di residenza umana
nonostante la singolare architettura techno-rétro, bassa e tondeggiante,
simile a una riproduzione in porcellana di una capanna di terra. Il
cielo alle spalle era di un turchese vivido, o almeno così lo aveva reso
la stampante. L'orizzonte era stranamente vicino ma geometricamente
piatto e diviso in rettangoli di verde coltivato che si diradavano,
coltivazioni che non riuscivo a riconoscere ma che erano troppo carnose
per essere grano o mais e troppo alte per essere lattuga o cavolo verza.
In primo piano c'erano due marziani adulti, maschio e femmina, con
un'espressione comicamente severa. Gotico marziano. Mancavano solo
forcone e firma di Grant Wood.

«Mia madre e mio padre» disse semplicemente Wun.

\emph{Foto 2}: «Io, da bambino.»

Questa era stupefacente. La pelle rugosa dei marziani, spiegò Wun, si
sviluppa in pubertà. Era Wun più o meno a sette anni, sorridente e con
la pelle liscia. Sembrava un bambino terrestre, anche se non era
possibile attribuirgli un'etnia - capelli biondi, pelle color caffè,
naso piccolo e labbra generose. Si trovava in quello che a prima vista
sembrava un eccentrico parco a tema, ma che in realtà era, secondo
quanto mi disse Wun, una città marziana. Un mercato. Bancarelle di cibo
e negozi, edifici di sgargianti colori primari fatti dello stesso
materiale simile alla porcellana della fattoria. La strada dietro di lui
era gremita di mezzi leggeri e pedoni. Fra gli edifici più alti era
visibile solo un lembo di cielo, e anche lì era stato colto il passaggio
di un qualche tipo di veicolo, con lame che giravano a giostra sfumate
in un pallido ovale.

«Sembri felice» dissi.

«La città si chiama Voy Voyud. Quel giorno arrivammo dalla campagna per
fare acquisti. Essendo primavera, i miei genitori mi permisero di
comprare dei \emph{murkud}. Animali di piccola taglia, da compagnia.
Sono simili alle rane. In quella borsa che reggo\ldots{} vedi?»

Wun stringeva un sacchetto di stoffa bianca che conteneva delle
misteriose protuberanze. \emph{Murkud}.

«Vivono solo qualche settimana, ma le loro uova sono deliziose.»

\emph{Foto 3}: Una vista panoramica. Nel terreno vicino, un'altra casa
marziana, una donna in un caffettano multicolore (sua moglie, mi spiegò)
e due ragazze carine dalla pelle liscia con abiti color ambra a forma di
sacco (le sue figlie.)

La fotografia era stata scattata dall'alto. Oltre la casa, era visibile
un intero paesaggio semirurale. Campi verdi e paludosi si crogiolavano
sotto un altro cielo turchese. La terra coltivata era divisa da
carreggiate elevate su cui viaggiavano pochi veicoli a forma di scatola,
e c'erano macchine agricole tra le colture, aggraziate mietitrici nere.
E all'orizzonte, dove le strade convergevano, c'era una città, la stessa
città, disse Wun, dove aveva comprato i \emph{murkud} da bambino: Voy
Voyud, la capitale della provincia di Kirioloj, con le sue alte torri a
terrazze intricate, che si ergevano nella bassa gravità.

«In questa foto si vede quasi tutto il delta del Kirioloj.» Il fiume era
una striscia azzurra che alimentava un lago del colore del cielo. La
città di Voy Voyud era stata costruita su un terreno elevato, il bordo
eroso di un antico cratere meteoritico, disse Wun, anche se a me
sembrava una linea ordinaria di colline basse.

I puntini neri sul lago lontano potevano essere barche, o chiatte.

«È un posto bellissimo» dissi.

«Sì.»

«Il paesaggio, ma anche la tua famiglia.»

«Sì.» Il suo sguardo incrociò il mio. «Sono morti.»

«Ah, mi dispiace.»

«Morirono in un'enorme alluvione diversi anni fa. L'ultima fotografia,
vedi? È la stessa vista, ma è stata scattata subito dopo il disastro.»

Una tempesta furiosa aveva scaricato piogge da record sulle pendici
delle Montagne Solitarie alla fine di una lunga stagione secca. Quasi
tutta quella pioggia era stata incanalata negli affluenti prosciugati
del Kirioloj. Marte terraformato, per certi aspetti, era ancora un mondo
giovane, stava ancora stabilizzando i suoi cicli idrologici e i suoi
paesaggi, che evolvevano rapidamente in polvere antica e regolite,
venivano riassettati dall'acqua in circolazione. Il risultato della
pioggia improvvisa ed eccessiva era stato un impasto di fango di ossido
rosso esondato dal Kirioloj che si era riversato nel delta agricolo come
un treno merci fluido.

\emph{Foto 4}: L'indomani. Della casa di Wun erano rimaste solo le
fondamenta e un unico muro, in piedi come cocci di ceramica in una
pianura caotica di fango, detriti e rocce. La lontana città sulla
collina era intatta, ma la terra fertile era stata sepolta. Fatta
eccezione per un luccichio di acqua marrone proveniente dal lago, quello
era Marte tornato quasi alla sua condizione vergine, un regolite senza
vita. Molti elicotteri sorvolavano la zona, presumibilmente alla ricerca
di sopravvissuti.

«Avevo trascorso una giornata sulle colline pedemontane con amici e
quando tornai a casa trovai tutto questo. Molte vite andarono perdute,
non solo la mia famiglia. E così, conservo queste quattro fotografie per
ricordarmi da dove vengo. Anche perché non posso tornare.»

«Dev'essere stato tremendo.»

«Me ne sono fatto una ragione. Per quel che è possibile. Quando ho
lasciato Marte, il delta era stato ripristinato. Non com'era un tempo,
ovviamente. Però era fertile, vivo, produttivo.»

Tutto lì, quello che sembrava volesse raccontare al riguardo.

Riguardai le immagini precedenti, rammentando a me stesso quello che
stavo osservando. Non erano effetti fantasiosi di immagini create al
computer, erano fotografie normali. Fotografie di un altro mondo. Di
Marte, un pianeta che per molto tempo avevamo gravato della nostra più
sconsiderata immaginazione. «Non è Burroughs, certamente non Wells,
forse un po' Bradbury\ldots»

Wun corrugò la fronte già fittamente aggrottata di suo. «Mi
spiace\ldots{} non conosco quelle parole.»

«Sono scrittori. Scrittori di romanzi, che hanno descritto il tuo
pianeta.»

Quando gli feci capire il concetto che certi autori avevano immaginato
la vita su Marte molto prima della sua effettiva terraformazione, Wun ne
rimase affascinato. «Sarebbe possibile leggere questi libri? E
discuterne la prossima volta che vieni a trovarmi?»

«Ne sarei lusingato. Sei sicuro di averne il tempo? Ci saranno
senz'altro capi di Stato desiderosi di parlarti.»

«Ne sono certo, ma possono aspettare.»

Gli dissi che non ne vedevo l'ora.

Tornando a casa, lungo la strada feci un salto in una libreria
dell'usato e al mattino consegnai un pacco di tascabili a Wun, o
perlomeno agli uomini taciturni che facevano la guardia al suo alloggio.

\emph{La guerra dei mondi, La principessa di Marte, Cronache marziane,
	Straniero in terra straniera, Il rosso di Marte}.

Per un paio di settimane non ebbi più sue notizie.

\separatorefiori

La costruzione delle nuove strutture alla Perielio proseguiva. Verso la
fine di settembre, al posto di una boscaglia di pini e palme nane
frastagliate c'erano massicce fondamenta in calcestruzzo, un'imponente
struttura di travi in acciaio e tubi in alluminio.

Molly aveva sentito che sarebbero arrivati per la settimana successiva
un laboratorio militare e attrezzature di refrigerazione. (Un'altra cena
al Champs, quasi tutti i clienti guardavano una partita dei Marlins
sullo schermo al plasma delle dimensioni di un tabellone pubblicitario
mentre noi ci dividevamo gli stuzzichini in un angolo lontano, buio.)
«Perché abbiamo bisogno di attrezzatura da laboratorio, Ty? La Perielio
si occupa di ricerca spaziale e dello Spin. Non capisco.»

«Non lo so. Nessuno ne parla.»

«Forse potresti chiederlo a Jason, uno dei pomeriggi che trascorrete
all'ala nord.»

Le avevo detto che mi stavo consultando con Jase, non che ero stato
adottato dall'ambasciatore marziano. «Non ho quel tipo di nullaosta per
le informazioni.» Né, ovviamente, l'aveva Molly.

«Comincio a pensare che non ti fidi di me.»

«Seguo solo le regole, Moll.»

«Giusto,» disse lei, «sei proprio un \emph{santo}\ldots»

\separatorefiori

Jason si fermò a casa mia senza preavviso, fortunatamente una sera che
Molly non c'era, per parlare della sua cura. Gli dissi quello che aveva
detto Malmstein, che non ci sarebbero stati problemi ad aumentare il
dosaggio ma avremmo dovuto fare attenzione agli effetti collaterali. La
malattia non si era arrestata e c'era un limite per reprimere i sintomi
oltre il quale non potevamo spingerci. Non significava che fosse
condannato, solo che prima o poi avrebbe dovuto condurre la sua vita in
maniera differente: accettare la malattia invece che sopprimerla. (Oltre
a questo, c'era un altro stadio di cui nessuno dei due aveva parlato: la
disabilità totale e la demenza.)

«Ho capito» disse Jason. Si sedette vicino alla finestra, dando
un'occhiata di tanto in tanto al suo riflesso sul vetro, con una gamba
lunga che penzolava dall'altra. «Mi serve solo qualche mese.»

«Qualche mese per cosa?»

«Qualche mese per tagliare le gambe a E.D. Lawton.» Lo fissai. Pensai
fosse uno scherzo ma non sorrideva. «Devo spiegartelo?»

«Se vuoi che ne capisca il senso, sì, spiegamelo.»

«Io e E.D. abbiamo opinioni divergenti sul futuro della Perielio. Per
lui la Perielio esiste per sostenere l'industria aerospaziale. Il morale
della favola è quello e lo è sempre stato. Non ha mai creduto che
potessimo fare qualcosa riguardo allo Spin.» Jason scrollò le spalle.
«Quasi certamente ha ragione lui, nel senso che non possiamo sistemare
la questione. Ma questo non significa che non possiamo capirla. Certo,
non possiamo combattere una guerra contro gli Ipotetici, almeno non nel
vero senso della parola, ma possiamo eseguire un po' di azioni
scientifiche di disturbo. Ecco il motivo della presenza di Wun.»

«Non ti seguo.»

«Wun non è solo un ambasciatore di benevolenza interplanetaria. È venuto
qui con una proposta di sforzo congiunto che potrebbe darci qualche
indizio sugli Ipotetici, per capire da dove vengono, cosa vogliono, cosa
stanno facendo a entrambi i pianeti. L'idea sta ottenendo pareri
contrastanti. E.D. cerca di affondarla, pensa che sia inutile e metta a
rischio quel che ci rimarrà del capitale politico dopo la
terraformazione.»

«Quindi sarai tu ad affondare lui?»

Jason sospirò. «Può sembrare crudele ma E.D. non capisce che il suo
tempo è passato. Mio padre rappresenta quello che serviva al mondo
vent'anni fa. E per questo lo ammiro. Ha realizzato cose stupefacenti,
incredibili. Senza E.D. a spronare i politici non ci sarebbe mai stata
una Perielio. Uno dei paradossi dello Spin è che alla lunga le
conseguenze del genio di E.D. Lawton gli si sono ritorte contro - se non
fosse mai esistito, non esisterebbe neanche Wun Ngo Wen. Non si tratta
di un mio conflitto edipico. So esattamente chi è mio padre e cos'ha
fatto. Si muove con disinvoltura nei corridoi del potere, gioca a golf
con Garland. Benissimo. Ma è anche un prigioniero. Un prigioniero della
sua stessa miopia. I suoi giorni da visionario sono finiti. Non approva
il progetto di Wun perché non si fida della tecnologia - non gli piace
nulla che non possa sottoporre a ingegneria inversa; non gli piace il
fatto che i marziani possano maneggiare tecnologie che noi solo ora
cominciamo a ipotizzare. E odia il fatto che io stia dalla parte di Wun.
Io e, potrei aggiungere, una nuova generazione di mediatori politici di
Washington, compreso Preston Lomax, che forse sarà il prossimo
Presidente. All'improvviso, E.D. è circondato da persone che non può
manipolare. Persone più giovani, persone che hanno assimilato lo Spin in
un modo che la sua generazione non ha mai fatto. Persone come noi, Ty.»

Ero un po' lusingato e un po' allarmato, a essere incluso in quel
pronome.

Gli dissi: «Te la stai prendendo troppo, mi pare\ldots»

Mi lanciò un'occhiata brusca. «Faccio esattamente quello che E.D. mi ha
addestrato a fare. Fin dalla nascita. Non ha mai voluto un figlio;
voleva un erede, un apprendista. Ha preso quella decisione molto tempo
prima dello Spin, Tyler. Sapeva esattamente quanto fossi brillante e
come dovesse sfruttare quell'intelligenza. E io mi sono adeguato. Anche
quando sono cresciuto abbastanza da capire in cosa era coinvolto, ho
collaborato. E così eccomi qui, una creazione di E.D. Lawton: l'oggetto
bello, esperto, asessuato, disponibile con i media che hai di fronte.
Ottimo prodotto di marketing, un certo acume intellettuale e nessuna
lealtà che non inizi e finisca con la Perielio. Ma c'è sempre stata una
piccola postilla su quel contratto, anche se a E.D. non piace pensarci.
``Erede'' implica ``eredità''. Implica che, a un certo punto, il mio
giudizio avrebbe rimpiazzato il suo. Be', è arrivato il momento.
L'opportunità che ci si prospetta è troppo importante per mandarla a
puttane.»

Notai che aveva le mani strette a pugno e le gambe gli tremavano, ma era
l'intensità dell'emozione o un sintomo della malattia? E ancora, quanta
parte di quel monologo era autentica e quanta il risultato dei
neurostimolanti che gli stavo prescrivendo?

«Sembri spaventato» disse Jason.

«Ma di quale tecnologia marziana stiamo parlando?»

Fece un gran sorriso. «È qualcosa di molto ingegnoso. Semibiologico. In
scala ridottissima. In sostanza, si tratta di retroazioni
autocatalitiche molecolari con programmazione contingente scritta nei
loro protocolli riproduttivi.»

«Jase, per piacere, parla come mangi.»

«Minuscoli replicatori artificiali.»

«Esseri viventi?»

«In un certo senso sì, esseri viventi. Esseri viventi artificiali che
possiamo lanciare nello spazio.»

«E quindi, cosa fanno, Jase?»

Fece un sorriso ancora più ampio. «Mangiano ghiaccio e cacano
informazioni.»